\section{Introduction}
\label{sec:intro}

A grocery recommender must choose a short list from a user's past purchases. Retaining more history can reveal recurring preferences, but it also allows more old products to compete for limited positions. An aggregate prediction score does not distinguish these effects. Moreover, users with longer observed histories tend to have larger historical catalogs, complicating comparisons between history-length groups.

Earlier experiments in this project established strong personal-history baselines. Time weighting helped prediction, whereas adding clustering and association-rule features did not reliably improve the full model. These findings motivated three research questions. First, how sensitive are history-group comparisons to differences in task composition? Second, what happens when scoring history and candidate eligibility change separately? Third, how much candidate reduction can a fixed short-list ranker tolerate, and where does the resulting performance claim stop applying?

We address these questions using a fixed-target evaluation. A policy may remove products from the eligible candidate pool, but it cannot remove those products from the true repeat-target set used to evaluate recommendations. A factorial experiment then changes the history used for scoring independently of the products allowed to compete. Finally, a validation-selected policy is frozen and tested on previously unused confirmation users. A subsequent, explicitly posthoc analysis examines individual changes and recommendation depth.

The contribution is empirical and methodological. We demonstrate that the effects of old candidates depend on scoring, quantify recovered and displaced true hits, and identify a candidate-reduction policy that retains average Top~10 recall within a declared tolerance. We also show why this result is narrow: most short lists remain unchanged, while excluded true targets become reachable in longer full-pool lists. We do not propose a new state-of-the-art algorithm or infer how shoppers psychologically forget products.

\subsection{Relation to prior work}
TIFU-KNN models personalized item frequency with temporal weighting and information from neighboring users~\citep{hu2020}. A next-basket recommendation reality check distinguishes repeat and exploratory recommendation and emphasizes the importance of task definition and evaluation~\citep{li2023}. These established ideas inform our task and baselines; we do not claim to introduce the repeat--explore distinction. We use LightGBM~\citep{ke2017} as a learned ranker and compare an independently written adaptation of TIFU-KNN in an earlier cohort. The present study focuses on separating ranking information from candidate competition and on explaining performance retention through exact hit accounting.
